\section{Composition and graph flattening}

\subsection{Logical composition by cut}

Suppose a fixed proof graph has permanent assumptions \(\Gamma\) and already
compiled node formulas \(\Delta=[D_1,\ldots,D_m]\).  A new node carries a local
object derivation
\[
  \Delta,\Gamma\objectvdash D_{m+1}.
\]
If every \(D_i\) is derivable from \(\Gamma\), simultaneous cut gives
\(\Gamma\objectvdash D_{m+1}\).

\begin{proposition}[Proof-node compilation]\label{prop:node-cut}
A reverse-topological sequence of local \(\Prf\,\PRAAx\) derivations compiles
to object derivations of every graph node from the permanent context.
\end{proposition}

\begin{proof}
Induct on the sequence.  The empty bundle is trivial.  At a node, apply
\leanname{Prf.trans} to replace each earlier-node hypothesis by its compiled
derivation and retain each permanent hypothesis by \leanname{Prf.hyp}.  Add the
new derivation to the bundle.
\end{proof}

The supplied \pathref{KernelCore.lean} implements exactly this proof.  Its
public constructors are \pathref{compileLocal} and
\pathref{compileAndExtend}; no custom graph logic is needed.

\subsection{Arithmetic flattening to a master certificate}

Logical cut presupposes that every node's finite arithmetic contract has been
implemented.  \(\FQNF\) offers a second, orthogonal normalization.

Consider a fixed acyclic graph with node certificates
\(c_i<B_i(q)\), after backward demand has converted all resource bounds to
query-only terms.  Pack the tuple \((c_0,\ldots,c_{m-1})\) by mixed radix.  The
global verifier:

\begin{enumerate}
\item decodes every \(c_i\);
\item checks \(c_i<B_i(q)\);
\item evaluates node \(i\)'s fixed checker using the decoded parent outputs;
\item checks the terminal specification.
\end{enumerate}

Because the graph is fixed and finite, these checks form one bounded Boolean
fold.  The global bound is \(\prod_iB_i(q)\).

\begin{theorem}[Master-certificate flattening]\label{thm:master-cert}
Let a fixed finite graph have query-only numerical bounds for every node and
reflected fixed local checkers.  Then the complete family of node certificates
has a single mixed-radix code with a reflected QF global checker.  If the graph
contracts have a simultaneous satisfying assignment, bounded search selects a
least accepted master code.
\end{theorem}

\begin{proof}
Mixed-radix coding gives the bound and projection equations.  Substitute the
projections into each local checker.  Closure of reflected checkers under
finite conjunction gives the global checker.  A satisfying assignment packs
to a witness below the product bound; apply \cref{thm:witness-erasure}.
\end{proof}

\subsection{Relationship between the two compositions}

The cut compiler and master-certificate compiler solve different problems.

\begin{itemize}
\item \textbf{Cut composition} is proof modularity.  Each node may use a
specialized algorithm and expose an arbitrary object formula.
\item \textbf{Master-code composition} is finite-data normalization.  It turns
all node evidence into one bounded witness and is especially useful before a
shared coding library is mature.
\end{itemize}

A practical implementation may combine them: use master codes inside
arithmetic-heavy regions, then expose concise node formulas and compose those
by cut.

\subsection{Backward demand}

For approximate or resource-sensitive nodes, a terminal query determines the
precision or bound required at each predecessor.  Since the graph is fixed,
this reverse pass is itself a fixed primitive-recursive computation.  It may
run in the host to generate specialized syntax, or it may be internalized if a
theorem needs the demand values as object terms.

\begin{proposition}[Demand soundness]
Suppose every node contract states that inputs meeting demands
\(d_1(p),\ldots,d_r(p)\) produce an output meeting demand \(p\).  The reverse
composition of these transformers assigns sufficient demands to every node of
a fixed DAG.
\end{proposition}

\begin{proof}
Induct in topological order during the forward proof.  The reverse pass has
already assigned each parent at least the demand requested by every child.
Apply the local contract at each node, then use cut.  No monotonicity is needed
unless multiple demands are merged by an inequality rather than computed
exactly; in that case the required monotonicity theorem is an explicit local
obligation.
\end{proof}

\subsection{Cycles and loops}

A genuine cycle is not a finite DAG node dependency.  It must be represented by
a bounded trace, a primitive recursion, or a fixed point with a proved finite
convergence bound.  Unrolling a query-bounded number of stages produces an
array certificate; direct recursion uses the next section's invariant theorem.
